\subsection{Polynomially convex sets}

DEFINITION 3. --- Let \(\Lambda\) a set and V a subset of \(\mathbf{C}^{\Lambda}\) . We say that V is polynomially convex if V is the set of points \((c_{\lambda})_{\lambda \in \Lambda}\) of \(\mathbf{C}^{\Lambda}\) such that

\[
| \mathrm{P} ((c _ {\lambda})) | \leqslant \sup _ {c \in \mathrm{V}} | \mathrm{P} (c) |
\]

for all \(\mathrm{P}\in \mathbf{C}[(\mathrm{X}_{\lambda})]\)

Lemma 3. --- Let \(\Lambda\) a set. A subset V of \(\mathbf{C}^{\Lambda}\) is polynomially convex if and only if there exist a family \((\mathrm{P}_i)_{i\in \mathrm{I}}\) of elements of \(\mathbf{C}[(\mathrm{X}_{\lambda})_{\lambda \in \Lambda}]\) and a family \((\mathrm{M}_i)_{i\in \mathrm{I}}\) of elements of \([0, +\infty]\) such that V is the set of \(c\in \mathbf{C}^{\Lambda}\) satisfying

\[
| \mathrm{P} _ {i} (c) | \leqslant \mathrm{M} _ {i}
\]

for all i ∈ I.

If the subset V of \(C^{\Lambda}\) is polynomially convex, it satisfies the above condition for the family consisting of the elements P of \(\mathbf{C}[(X_{\lambda})_{\lambda\in\Lambda}]\) by setting \(M_{P}=\sup_{c\in V}|P(c)|\) .

Conversely, let \((\mathrm{P}_{i})_{i\in\mathrm{I}}\) be a family of elements of \(\mathbf{C}[(\mathrm{X}_{\lambda})_{\lambda\in\Lambda}]\) and \((\mathrm{M}_{i})_{i\in\mathrm{I}}\) be a family of elements of \([0,+\infty]\) . Let V be the set of c in \(C^{\Lambda}\) such that \(|\mathrm{P}_{i}(c)|\leqslant\mathrm{M}_{i}\) for all \(i\in I\) . We then have \(\sup_{c\in\mathrm{V}}|\mathrm{P}_{i}(c)|\leqslant\mathrm{M}_{i}\) for \(i\in I\) . Suppose that \(x\in C^{\Lambda}\) satisfies

\[
| \mathrm{P} (x) | \leqslant \sup _ {c \in \mathrm{V}} | \mathrm{P} (c) |
\]

for all \(\mathrm{P}\in\mathbf{C}[(\mathrm{X}_{\lambda})]\) . For \(i\in I\) , we have in particular \(|\mathrm{P}_{i}(x)|\leqslant\mathrm{M}_{i}\) , hence \(x\in V\) . Conversely, for every element \(x\in V\) and every polynomial \(\mathrm{P}\in\mathbf{C}[(\mathrm{X}_{\lambda})]\) , one has \(|\mathrm{P}(x)|\leqslant\sup_{c\in\mathrm{V}}|\mathrm{P}(c)|\) Consequently, the set V is polynomially convex.

Lemma 4. --- Let A be a unital commutative Banach algebra. Let \(\Lambda\) a set and \(x = (x_{\lambda})_{\lambda \in \Lambda}\) a family of elements of A. If the unital subalgebra generated by the family x is dense in A, then the simultaneous spectrum \(\mathrm{Sp}_{\mathrm{A}}^{\Lambda}(x)\) is polynomially convex.

This follows from the assertion. \(b\) ) of prop. 12 of I, p.~43 and of definition 3.

Any intersection of polynomially convex subsets of \(C^{\Lambda}\) is polynomially convex (Lemma 3). This justifies the following definition:

DEFINITION 4. --- Let \(\Lambda\) a set and V a subset of \(\mathbf{C}^{\Lambda}\) The polynomially convex hull of V is the smallest polynomially convex subset of \(\mathbf{C}^{\Lambda}\) containing V.

The polynomially convex hull of V is the set of c belonging to \(C^{\Lambda}\) such that \(|P(c)| \leqslant \sup_{W} |P|\) for all \(P \in \mathbf{C}[(X_{\lambda})]\) Indeed, this set is polynomially convex by Lemma 3, and is contained in every polynomially convex set containing V by definition.

Example. --- Let \(\Lambda\) a finite set and \(V \subset \mathbf{C}^{\Lambda}\) a compact convex subset. Then \(V\) is polynomially convex. Indeed, let \(W\) be the polynomially convex hull of \(V\) . Let us prove that \(W \subset V\) , which will imply the assertion. Let \(x \in \mathbf{C}^{\Lambda} - V\) . There exists a real hyperplane \(H\) into \(\mathbf{C}^{\Lambda}\) which strictly separates \(x\) and \(V\) (EVT, II, p.~41, prop. 4). Let \(f_{\mathbf{R}}\) a form \(\mathbf{R}\) -linear on \(\mathbf{C}^{\Lambda}\) and \(\alpha \in \mathbf{R}\) such that \(H\) is the set of \(y \in \mathbf{C}^{\Lambda}\) satisfying \(f_{\mathbf{R}}(y) = \alpha\) . Let \(f\) a linear form on \(\mathbf{C}^{\Lambda}\) such that \(f_{\mathbf{R}} = \mathcal{R}(f)\) . We therefore have

\[
\mathcal {R} (f (x)) > \sup _ {y \in V} \mathcal {R} (f (y)).
\]

For every \(t \in R\) and \(y \in V\) , set \(f_{t}(y) = t + f(y)\) . We have \(|f_{t}| - \mathcal{R}(f_{t}) \to 0\) into \(\mathcal{C}(\mathbf{C}^{\Lambda}, \mathbf{R})\) endowed with the topology of compact convergence when \(t \to +\infty\) . For t sufficiently large, it follows that \(|f_{t}(x)| > \sup_{y \in V} |f_{t}(y)|\) since V is compact. Thus \(x \in C^{\Lambda} - W\) since \(f_{t}\) is a polynomial function.

Lemma 5. --- Let K be a compact subset of C. We denote by \(\widehat{\mathbf{K}}\) the union of K and the connected components of \(\mathbf{C} - \mathbf{K}\) which are relatively compact. Then the set \(\widehat{\mathbf{K}}\) is compact.

Since K is compact, there exists a real number r > 0 such that K is contained in the open disk D centered at 0 with radius r. Then C-K contains C-D. Since the space C-D is connected (being homeomorphic to \([r, +\infty[\times\mathbf{S}^1)\) , it is contained in a connected component U of \(\mathbf{C}-\mathbf{K}\) . Every other connected component of \(\mathbf{C}-\mathbf{K}\) is contained in D, hence is bounded. The connected component U is then the unique unbounded connected component of \(\mathbf{C}-\mathbf{K}\) , that is, \(\mathrm{U} = \mathbf{C}-\widehat{\mathrm{K}}\) . Since U is open and contains the complement of the disk D, the set \(\widehat{\mathrm{K}}\) is compact.

PROPOSITION 13. --- Let \(n \geqslant 1\) be an integer. Let \(\mathbf{K} \subset \mathbf{C}^{n}\) be a closed subset and V its polynomially convex hull.

\begin{enumerate}
\def\labelenumi{\alph{enumi})}
\item
  Every bounded connected component of \(C^{n} - K\) is contained in V;
\item
  If n = 1, and if K is compact, then V is the union \(\widehat{K}\) of K and the bounded connected components of C-K.
\end{enumerate}

Since \(K \subset C^{n}\) is closed, its complement \(C^{n}-K\) is open, hence locally connected, so that each connected component of \(C^{n}-K\) is open. The maximum principle (VAR, R1, p.~29, 3.3.7) then shows that every bounded connected component of \(C^{n}-K\) is contained in V, which proves assertion a).

Now suppose \(n = 1\) and K compact. Let \(\widehat{\mathbf{K}}\) denote the union of K and the bounded connected components of \(\mathbf{C} - \mathbf{K}\) , so that \(\widehat{\mathbf{K}} \subset \mathbf{V}\) ; by the preceding, the set \(\widehat{\mathbf{K}}\) is compact (lemma 5). Let A be the commutative unital Banach algebra \(\mathcal{C}(\widehat{\mathbf{K}})\) . Let \(x \in \mathbf{A}\) the identity function on \(\widehat{\mathbf{K}}\) and B the closed unital subalgebra of A generated by \(x\) . We have \(\mathrm{Sp}_{\mathrm{A}}(x) = \widehat{\mathbf{K}}\) (example 3 of I, p.~17), hence \(\mathbf{C} - \mathrm{Sp}_{\mathrm{A}}(x)\) is connected. Consequently we have \(\mathrm{Sp}_{\mathrm{B}}(x) = \mathrm{Sp}_{\mathrm{A}}(x)\) (cor. of prop. 6 of I, p.~28). Since \(\mathrm{Sp}_{\mathrm{B}}(x)\) is polynomially convex by lemma 4, this shows that \(\widehat{\mathbf{K}}\) is polynomially convex and therefore that \(\mathbf{V} \subset \widehat{\mathbf{K}}\) .

The second part of the proposition does not extend to the case \(n \geqslant 2\) (cf.~exerc. 23 of I, p.~170).

PROPOSITION 14. --- Let \(\Lambda\) be a set. Let A be a commutative unital Banach algebra, \(x = (x_{\lambda})_{\lambda \in \Lambda}\) a family of elements of A and A' the unital Banach subalgebra generated by \(x\) . Then the simultaneous spectrum \(\mathrm{Sp}_{\mathrm{A}'}^{\Lambda}(x)\) is the polynomially convex hull of \(\mathrm{Sp}_{\mathrm{A}}^{\Lambda}(x)\) .

Indeed, prop. 12 of I, p.~43, b) proves that \(\mathrm{Sp}_{\mathrm{A}'}^{\Lambda}(x)\) is the set of \(c \in \mathbf{C}^{\Lambda}\) such that \(|\mathrm{P}(c)| \leqslant \varrho(\mathrm{P}(x))\) for all \(\mathrm{P} \in \mathbf{C}[(\mathrm{X}_{\lambda})]\) .

Now we have

\[
\varrho (\mathrm{P} (x)) = \sup _ {\chi \in \mathsf {X} (\mathrm{A})} | \chi (\mathrm{P} (x)) | = \sup _ {\chi \in \mathsf {X} (\mathrm{A})} | \mathrm{P} ((\chi (x _ {\lambda})) _ {\lambda \in \Lambda}) | = \sup _ {c \in \operatorname{Sp} _ {\mathrm{A}} ^ {\Lambda} (x)} | \mathrm{P} (c) |,
\]

by prop. 7 of I, p.~38, a) and prop. 12 of I, p.~43, a). The result then follows from lemma 3.

PROPOSITION 15. --- Let \(\Lambda\) be a set. Let K be a compact polynomially convex subset of \(\mathbf{C}^{\Lambda}\) . Let \(\mathrm{A}_0\) the set of restrictions to K of polynomial functions on \(\mathbf{C}^{\Lambda}\) , and let A be the closure of \(\mathrm{A}_0\) in algebra \(\mathcal{C}(\mathrm{K})\) .

Let ev: K → X(A) be the map defined by x → ev\_x, where ev\_x is the character f → f(x) of A. Let z = (z\_λ)\_\{λ∈Λ\} be the family in A of restrictions to K of the coordinate functions on C\^{}Λ, and let φ be the surjective map from X(A) to Sp\_A\^{}Λ(z) defined by the family z.

Then we have K = Sp \(_{A}^{\Lambda}(z)\) , and the maps ev: K → X(A) and φ: X(A) → K are reciprocal homeomorphisms.

The map \(\varphi \circ\) ev is the identity map of K. In particular, K is contained in the image \(\mathrm{Sp}_{\mathrm{A}}^{\Lambda}(z)\) of \(\varphi\) .

Prop. 14 implies that \(\mathrm{Sp}_{\mathrm{A}}^{\Lambda}(z)\) is the polynomially convex hull of \(\mathrm{Sp}_{\mathcal{C}(\mathrm{K})}^{\Lambda}(z)\) . Since, on the other hand, we have \(\mathrm{Sp}_{\mathcal{C}(\mathrm{K})}^{\Lambda}(z) = \mathrm{K}\) (I, p.~42, example), which is polynomially convex by hypothesis, we deduce that \(\mathrm{Sp}_{\mathrm{A}}^{\Lambda}(z) = \mathrm{K}\) .

Since the unital subalgebra generated by the family of elements \(z_{\lambda}\) is dense in A, the map \(\varphi\) is a homeomorphism from \(\mathsf{X}(\mathsf{A})\) on \(\mathrm{Sp}_{\mathsf{A}}^{\Lambda}(z)=\mathrm{K}\) Proposition 12 of Book I, p. 43(a). The identity \(\varphi\circ ev=Id_{K}\) then shows that ev is the inverse homeomorphism of \(\varphi\) .

For every subset \(\Lambda'\) of \(\Lambda\) , we denote by \(\mathrm{pr}_{\Lambda'}\) the canonical projection \(\mathbf{C}^{\Lambda} \to \mathbf{C}^{\Lambda'}\) . Let W be a subset of \(\mathbf{C}^{\Lambda}\) and let V be its polynomially convex hull. Set \(W' = \mathrm{pr}_{\Lambda'} W\) . Since every element of \(\mathbf{C}[(X_{\lambda})_{\lambda \in \Lambda'}]\) is identified with an element of \(\mathbf{C}[(X_{\lambda})_{\lambda \in \Lambda}]\) the polynomially convex hull of \(W'\) is contained in \(\mathrm{pr}_{\Lambda'} V\) .

Lemma 6. --- Let \(K \subset C^{\Lambda}\) a compact polynomially convex subset and U a neighborhood of K. There exists a finite subset \(\Lambda_{0}\) of \(\Lambda\) such that, for every subset \(\Lambda'\) of \(\Lambda\) containing \(\Lambda_{0}\) , the set \(\mathrm{pr}_{\Lambda'}(\mathrm{U})\) contains the polynomially convex hull of \(\mathrm{pr}_{\Lambda'}(\mathrm{K})\) .

Since K is compact, there exists a family of compact disks \(D_{\lambda}\) in C with center 0 and radii \(R_{\lambda}\) such that K is contained in the product \(\mathrm{D} = \prod_{\lambda} \mathrm{D}_{\lambda}\) . For every \(\mathrm{P} \in \mathbf{C}[(\mathrm{X}_{\lambda})]\) , let \(\mathrm{K}_{\mathrm{P}}\) the set of \(x \in \mathbf{C}^{\Lambda}\) such that

\[
| \mathrm{P} (x) | \leqslant \sup _ {c \in \mathrm{K}} | \mathrm{P} (c) |.
\]

We have

\[
\mathrm{D} \cap \bigcap_ {\mathrm{P}} \mathrm{K} _ {\mathrm{P}} = \mathrm{K}\tag{3}
\]

which is contained in U. Thus, the space D is the union of the open set \(D \cap U\) and of the family of open sets \(D - K_{P}\) for \(P \in \mathbf{C}[(X_{\lambda})]\) Since D is compact, there exists an integer \(q \geqslant 1\) and polynomials \(P_{1}, \ldots, P_{q} \in \mathbf{C}[(X_{\lambda})]\) such that:

\[
\mathrm{D} \cap \mathrm{K} _ {\mathrm{P} _ {1}} \cap \dots \cap \mathrm{K} _ {\mathrm{P} _ {q}} \subset \mathrm{U}.\tag{4}
\]

There exists a finite set \(\Lambda_0 \subset \Lambda\) such that \(P_i \in \mathbf{C}[(X_\lambda)_{\lambda \in \Lambda_0}]\) for \(1 \leqslant i \leqslant q\) . Let us show that \(\Lambda_0\) has the required property.

Let \(\Lambda'\) a subset of \(\Lambda\) containing \(\Lambda_0\) . Let E be the subset of \(\mathbf{C}^{\Lambda'}\) consisting of the elements \(c = (c_{\lambda})_{\lambda \in \Lambda'}\) defined by the inequalities \(|c_{\lambda}| \leqslant \mathrm{R}_{\lambda}\) for \(\lambda \in \Lambda'\) , and

\[
| \mathrm{P} _ {i} (c) | \leqslant \sup _ {c \in \mathrm{K}} | \mathrm{P} _ {i} (c) |
\]

for \(i = 1, \ldots, q\) . The set E is polynomially convex (lemma 3), and formula (3) shows that \(\mathrm{pr}_{\Lambda'}(\mathrm{K}) \subset \mathrm{E}\) .

On the other hand, let \(c = (c_{\lambda})_{\lambda \in \Lambda'} \in \mathrm{E}\) ; let \(d = (d_{\lambda})_{\lambda \in \Lambda}\) be the element of \(\mathbf{C}^{\Lambda}\) defined by \(d_{\lambda} = c_{\lambda}\) for \(\lambda \in \Lambda'\) and \(d_{\lambda} = 0\) for \(\lambda \in \Lambda - \Lambda'\) Then (4) implies that \(d \in \mathrm{U}\) and therefore that \(c \in \mathrm{pr}_{\Lambda'}(\mathrm{U})\) Thus, \(\mathrm{E} \subset \mathrm{pr}_{\Lambda'}(\mathrm{U})\) , which completes the proof.

Lemma 7. --- Let \(n \geqslant 1\) an integer and K a polynomially convex compact subset of \(\mathbf{C}^{n}\) Then K admits a fundamental system of compact neighborhoods that are polynomially convex.

There exists a compact polydisk (cf.~VAR, R1, p.~24) \(\Delta\) of \(\mathbf{C}^n\) which is a neighborhood of K. Since K is polynomially convex, there exists a family \((\mathrm{P}_i)_{i\in \mathrm{I}}\) of elements of \(\mathbf{C}[\mathrm{X}_1,\dots ,\mathrm{X}_n]\) , and a family \((\mathrm{M}_i)_{i\in \mathrm{I}}\) of positive real numbers such that K is the set of \(z\in\) \(\Delta\) satisfying \(|\mathrm{P}_i(z)|\leqslant \mathrm{M}_i\) for all \(i\) (Lemma 3). For every finite subset J of I and every \(\varepsilon >0\) , let \(\mathrm{K}_{\mathrm{J},\varepsilon}\) the set of \(z\in \Delta\) such that \(|\mathrm{P}_i(z)|\leqslant \mathrm{M}_i + \varepsilon\) for \(i\in \mathrm{J}\) . Then each set \(\mathrm{K}_{\mathrm{J},\varepsilon}\) is a compact polynomially convex neighborhood of K (loc. cit.), and the intersection of the sets \(\mathrm{K}_{\mathrm{J},\varepsilon}\) is K. The sets \(\mathrm{K}_{\mathrm{J},\varepsilon}\) therefore form a fundamental system of polynomially convex neighborhoods of K (TG, I, p.~60, th. 1).

